\documentclass[answers]{exam}
\usepackage{../../../mypackages}
\usepackage{../../../macros}

\title{Correction --- Interrogation N°4 (Calcul mental)}
\author{N. Bancel}
\date{16 Octobre 2025}

% Mise en forme des solutions en bleu
\SolutionEmphasis{\color{blue}}
\renewcommand{\solutiontitle}{\noindent}

\begin{document}

\textbf{Collège Lycée Suger}
\hfill
\textbf{Mathématiques} \\

\textbf{Année 2025-2026}
\hfill
\textbf{Classe de 6ème} \par

{\let\newpage\relax\maketitle}

\begin{center}
  \textbf{\textcolor{red}{Correction de l'interrogation N°4 de calcul mental}} \\
  Les solutions en \textcolor{blue}{bleu} expliquent une méthode possible pour trouver le résultat.
\end{center}

\section*{Correction détaillée}

\begin{questions}

  \question $41 \times 4 = \boxed{164}$

  \begin{solution}
  \begin{align*}
    41 \times 4 
      &= 41 \times (2 \times 2) \\
      &= (41 \times 2) \times 2 &&\text{on remplace $4$ par $2 \times 2$} \\
      &= 82 \times 2 &&\text{car } 41 \times 2 = 82 \\
      &= 164.
  \end{align*}
  \end{solution}

  \question $2{,}5 \times 6{,}68 \times 4 = \boxed{66{,}8}$

  \begin{solution}
  \begin{align*}
    2{,}5 \times 6{,}68 \times 4
      &= 2{,}5 \times 6{,}68 \times (2 \times 2) \\
      &= (2{,}5 \times 2) \times 6{,}68 \times 2 &&\text{on regroupe un premier facteur $2$ avec $2{,}5$} \\
      &= 5 \times 6{,}68 \times 2 &&\text{car } 2{,}5 \times 2 = 5 \\
      &= (5 \times 2) \times 6{,}68 &&\text{on regroupe le deuxième facteur $2$ avec $5$} \\
      &= 10 \times 6{,}68 \\
      &= 66{,}8 &&\text{multiplier par $10$, c'est décaler la virgule d'un rang vers la droite.}
  \end{align*}
  \end{solution}

  \question $284 \div 4 = \boxed{71}$

  \begin{solution}
  On utilise la méthode « diviser deux fois par $2$ » :
  \begin{align*}
    284 \div 4 
      &= 284 \div (2 \times 2) \\
      &= (284 \div 2) \div 2 &&\text{on partage d'abord en $2$ parts égales, puis encore en $2$} \\
      &= 142 \div 2 &&\text{car } 284 \div 2 = 142 \\
      &= 71.
  \end{align*}
  \end{solution}

  \question $21 + 3{,}1 + 79 + 0{,}9 = \boxed{104}$

  \begin{solution}
  On regroupe les entiers et les décimaux séparément :
  \begin{align*}
    21 + 3{,}1 + 79 + 0{,}9
      &= (21 + 79) + (3{,}1 + 0{,}9) \\
      &= 100 + 4 &&\text{car } 21 + 79 = 100 \text{ et } 3{,}1 + 0{,}9 = 4 \\
      &= 104.
  \end{align*}
  \end{solution}

  \question $2{,}58 \times 1\,000 = \boxed{2\,580}$

  \begin{solution}
  Multiplier par $1\,000$, c'est déplacer la virgule de \textbf{trois rangs vers la droite} :
  \begin{align*}
    2{,}58 \times 1\,000 &= 2\,580.
  \end{align*}
  On peut aussi voir $1\,000 = 10^3$ : on gagne trois zéros.
  \end{solution}

  \question $4{,}06 \times 100 = \boxed{406}$

  \begin{solution}
  Multiplier par $100$, c'est déplacer la virgule de \textbf{deux rangs vers la droite} :
  \begin{align*}
    4{,}06 \times 100 &= 406.
  \end{align*}
  On peut le voir en deux étapes :
  \[
    4{,}06 \times 10 = 40{,}6 \quad\text{puis}\quad 40{,}6 \times 10 = 406.
  \]
  \end{solution}

  \question $96 \div 4 = \boxed{24}$

  \begin{solution}
  On divise deux fois par $2$ :
  \begin{align*}
    96 \div 4 
      &= 96 \div (2 \times 2) \\
      &= (96 \div 2) \div 2 \\
      &= 48 \div 2 &&\text{car } 96 \div 2 = 48 \\
      &= 24.
  \end{align*}
  \end{solution}

  \question $25 + (7 - (3 + 2)) = \boxed{27}$

  \begin{solution}
  On suit les priorités de calcul avec un alignement ligne par ligne :
  \begin{align*}
    25 + (7 - (3 + 2))
      &= 25 + (7 - 5) &&\text{car } 3 + 2 = 5 \\
      &= 25 + 2 &&\text{car } 7 - 5 = 2 \\
      &= 27.
  \end{align*}
  \end{solution}

  \question $897 \div 1\,000 = \boxed{0{,}897}$

  \begin{solution}
  Diviser par $1\,000$, c'est déplacer la virgule de \textbf{trois rangs vers la gauche}. \\
  On peut écrire $897$ sous la forme $897{,}0$ :
  \begin{align*}
    897{,}0 \div 1\,000 &= 0{,}897.
  \end{align*}
  \end{solution}

  \question $3 \times (10 - 2) + (7 - 4) = \boxed{27}$

  \begin{solution}
  On utilise un enchaînement de calculs avec parenthèses :
  \begin{align*}
    3 \times (10 - 2) + (7 - 4)
      &= 3 \times 8 + (7 - 4) &&\text{car } 10 - 2 = 8 \\
      &= 3 \times 8 + 3 &&\text{car } 7 - 4 = 3 \\
      &= 24 + 3 &&\text{car } 3 \times 8 = 24 \\
      &= 27.
  \end{align*}
  \end{solution}

\end{questions}

\end{document}
